% Manual for the mareimekomot package. Compile with XeLaTeX (twice).
\documentclass[11pt]{article}
\usepackage{fontspec}
\usepackage{polyglossia}
\setdefaultlanguage{english}
\setotherlanguage{hebrew}
\setmainfont{FreeSerif}
\newfontfamily{\hebrewfont}{FreeSerif}[Script=Hebrew]
\setmonofont{FreeMono}
\newfontfamily{\hebrewfonttt}{FreeMono}[Script=Hebrew]
\usepackage{mareimekomot}
\usepackage[margin=1in]{geometry}
\usepackage{fancyvrb}
\usepackage{booktabs}
\usepackage{array}
\usepackage{longtable}
\usepackage[hidelinks]{hyperref}
\fvset{fontsize=\small, frame=single, framesep=2mm}
\newcommand{\cs}[1]{\texttt{\char`\\#1}}
% long command names don't hyphenate
\setlength{\emergencystretch}{3em}

\ExplSyntaxOn
% \row{code}: the code, then what it gives in English and in Hebrew
\str_new:N \l__doc_code_str
\NewDocumentCommand \row { m }
	{ \__doc_code:n {#1} & #1 & \texthebrew {#1} \\ }
\cs_new_protected:Npn \__doc_code:n #1
	{
		\str_set:Nn \l__doc_code_str {#1}
		% \detokenize's space after a command name
		\regex_replace_all:nnN { (\\[A-Za-z]+) \x20 } { \1 } \l__doc_code_str
		\texttt { \l__doc_code_str }
	}
% \namelist{book}: a table of the names mareimekomot.sty declares with
% \__mm_book:nnn {keys} {Hebrew} {English} (or, for months, with a number
% after them), read from the package itself
\ior_new:N \g__doc_ior
\seq_new:N \l__doc_match_seq
\tl_new:N \l__doc_rows_tl
\tl_new:N \l__doc_he_tl
\tl_new:N \l__doc_en_tl
\cs_generate_variant:Nn \tl_set_rescan:Nnn { Nne }
\NewDocumentCommand \namelist { m }
	{
		\tl_clear:N \l__doc_rows_tl
		\ior_open:Nn \g__doc_ior { mareimekomot.sty }
		\ior_str_map_inline:Nn \g__doc_ior
			{
				\regex_extract_once:nnNT
					{ \A \s* \\__mm_ #1 :nnnn? \s* \{ ([^\}\#]*) \} \s* \{ (.*?) \} \s* \{ (.*?) \} (?: \s* \{ \s* [0-9]+ \s* \} )? \s* \Z }
					{##1} \l__doc_match_seq
					{
						\tl_set_rescan:Nne \l__doc_he_tl { \ExplSyntaxOn } { \seq_item:Nn \l__doc_match_seq { 3 } }
						\tl_set_rescan:Nne \l__doc_en_tl { \ExplSyntaxOn } { \seq_item:Nn \l__doc_match_seq { 4 } }
						\tl_put_right:Ne \l__doc_rows_tl
							{
								\exp_not:N \texttt { \tl_trim_spaces:e { \seq_item:Nn \l__doc_match_seq { 2 } } } &
								\exp_not:V \l__doc_en_tl &
								\exp_not:N \texthebrew { \exp_not:V \l__doc_he_tl } \exp_not:N \\
							}
					}
			}
		\ior_close:N \g__doc_ior
		\begin{longtable}
			{ @{} >{\raggedright\arraybackslash}p{0.38\linewidth}
				>{\raggedright\arraybackslash}p{0.36\linewidth}
				>{\raggedright\arraybackslash}p{0.20\linewidth} @{} }
		\toprule
		Keys & English & Hebrew \\
		\midrule
		\endhead
		\l__doc_rows_tl
		\bottomrule
		\end{longtable}
	}
\ExplSyntaxOff

\newenvironment{examples}
	{\begin{center}\small\begin{tabular}{@{}lll@{}}\toprule Command & English & Hebrew \\ \midrule}
	{\bottomrule\end{tabular}\end{center}}

\title{The \texttt{mareimekomot} package\\[0.3em]\large References to rabbinic texts, in Hebrew and English}
\author{Gabriel Gendler Yom-Tov\\\small\texttt{shiurtex@thm.mozmail.com}}
\date{Version 1.1, 29 October 2026}

\begin{document}

\maketitle
\tableofcontents

\section{Introduction}

\texttt{mareimekomot} (\texthebrew{מראי מקומות}, ``references'') writes references to Tanakh,
Talmud, halakhah, Midrash and later works, and the authors and dates that go with them. A
reference is written once and comes out in the language of the text around it:
\begin{examples}
\row{\tanach{gen}{37:18-20}}
\row{\bavli{menachot}{69b}}
\row{\shulchanaruch{oh}{605:1}}
\row{\sefer{bereshitrabbah}{38:13}}
\end{examples}
So a book or paper can cite the same way in its Hebrew and English passages, and a bilingual
source sheet can give each source one heading that serves both languages. The package was
written for the \texttt{shiur} package, which typesets source sheets and loads it, but it needs
nothing from \texttt{shiur} and works in any document.

The printed English follows one transliteration: \textit{ḥ} for \texthebrew{ח} and
\textit{kh} for \texthebrew{כ} (Menaḥot, Berakhot), and \textit{e} for a vocal shva (Seforno,
Keli Yakar). The keys you type are matched loosely, so any usual spelling works
(section~\ref{sec:keys}).

\subsection{Loading}

The package needs XeLaTeX or LuaLaTeX, and \texttt{polyglossia} with Hebrew as one of the
languages, and a font with Hebrew:
\begin{Verbatim}
\usepackage{polyglossia}
\setdefaultlanguage{english}
\setotherlanguage{hebrew}
\newfontfamily{\hebrewfont}{FreeSerif}[Script=Hebrew]
\usepackage{mareimekomot}
\end{Verbatim}
A reference is Hebrew wherever the current language is \texttt{hebrew}: inside
\verb|\texthebrew{...}| or a \texttt{hebrew} environment, or throughout a document whose main
language is Hebrew. Everywhere else it is English. Without Hebrew loaded, everything comes out
in English.

\subsection{Options}

Three settings change the style. They can be given as package options or, at any point in a
document, with \cs{mareimekomotsetup}:
\begin{Verbatim}
\usepackage[sham=full, daf=short]{mareimekomot}
\mareimekomotsetup{sham=full}
\end{Verbatim}
\begin{description}
\item[\texttt{sham=short} or \texttt{full}] Whether \cs{sham} (``ad loc'') also gives the
reference it stands for, as ``ad loc (\bavli{menachot}{69b})'' (section~\ref{sec:sham}). The
default is \texttt{short}.
\item[\texttt{daf=long} or \texttt{short}] Whether a daf in Hebrew is
\texthebrew{\bavli*{69b}} (\texttt{long}, the default) or the traditional
\texthebrew{\mareimekomotsetup{daf=short}\bavli*{69b}} with a colon for the second side and a
full stop for the first.
\item[\texttt{attribution=auto} or \texttt{off}] Whether a package that typesets sources, such
as \texttt{shiur}, gives the author of a work cited in a source's heading under the heading
(section~\ref{sec:authors}). The default is \texttt{auto}.
\end{description}

\section{Keys and spelling}\label{sec:keys}

Books, tractates, sections, commentaries and works are named by keys, written in English
letters: \texttt{gen}, \texttt{menachot}, \texttt{oh}, \texttt{rashi}, \texttt{tanchuma}. The
full lists are in the appendix. Keys are matched loosely, so that whatever transliteration
you use works:
\begin{itemize}
\item case, spaces, hyphens, full stops and apostrophes don't matter: \texttt{Bava Metzia},
	\texttt{bava-metzia}, \texttt{bavametzia};
\item \textit{ch}, \textit{kh} and \textit{h} are the same, as are \textit{tz}, \textit{ts}
	and \textit{z}, and \textit{y} and \textit{i}: \texttt{menachot}, \texttt{menakhot},
	\texttt{menahot};
\item doubled letters count as single, and a final \textit{h} doesn't matter:
	\texttt{sukkah}, \texttt{sukka}, \texttt{suka};
\item a numeral at either end can be a number or Roman: \texttt{samuelii}, \texttt{2sam},
	\texttt{samuel2}.
\end{itemize}
Some names have more keys, for forms that differ by more than that: \texttt{berachot} and
\texttt{brachot}, \texttt{tosafot} and \texttt{tosfot}. A key the package doesn't know is an
error, which names it.

Commands whose names have \texthebrew{ח} or \texthebrew{כ} can be spelt with \textit{ch} or
\textit{kh}: \cs{tanach} and \cs{tanakh}, \cs{shulchanaruch} and \cs{shulkhanarukh},
\cs{aruchhashulchan} and \cs{arukhhashulkhan}, \cs{shach} and \cs{shakh}, and \cs{bach},
\cs{bakh} and \cs{bah}.

\section{Tanakh}

\cs{tanach}\verb|{book}{reference}| takes a chapter (\texttt{14}), a verse (\texttt{37:18}), or
a range within a chapter (\texttt{37:18-20}) or across chapters (\texttt{37:18-38:2}, or
\texttt{37-38}). The starred form leaves out the book, for an addition such as
``(+\tanach*{gen}{39:1})'':
\begin{examples}
\row{\tanach{gen}{37:18-20}}
\row{\tanach{exodus}{14}}
\row{\tanach{kings ii}{4:1-5:3}}
\row{\tanach*{gen}{39:1}}
\end{examples}
A commentary on a verse goes in the optional argument (section~\ref{sec:commentaries}):
\begin{examples}
\row{\tanach[rashi]{gen}{18:7}}
\row{\tanach[onkelos]{deut}{6:5}}
\end{examples}

\section{Talmud and Mishnah}

\cs{bavli}\verb|{tractate}{daf}| takes a daf as \texttt{69b}, \texttt{69}, \texttt{69b-70a} or
\texttt{9a-b}, and \cs{bavli*}\verb|{daf}| gives the daf alone. \cs{mishnah} and \cs{yerushalmi}
take a chapter and its mishnah or halakhah, written like a verse:
\begin{examples}
\row{\bavli{menachot}{69b}}
\row{\bavli{yoma}{9a-b}}
\row{\bavli*{127b}}
\row{\mishnah{sukkah}{4:1}}
\row{\yerushalmi{berachot}{1:1-3}}
\end{examples}
A commentary on a daf goes in \cs{bavli}'s optional argument. In Hebrew, a commentary is on the
Bavli anyway, so \texthebrew{בבלי} is left out:
\begin{examples}
\row{\bavli[tosafot]{menachot}{69b}}
\row{\bavli*[rashi]{70a}}
\end{examples}

\section{Codes}

\cs{rambam}\verb|{laws}{reference}| takes a set of laws of the Mishneh Torah, by a key such as
\texttt{shofar}, \texttt{teshuvah} or \texttt{yesodeihatorah} (all 83 are in the appendix), and a
chapter and law:
\begin{center}
\verb|\rambam{shofar}{2:6}|\\[0.3em]
\rambam{shofar}{2:6}\\
\texthebrew{\rambam{shofar}{2:6}}
\end{center}

The Tur, the Shulḥan Arukh and the works arranged by them take a section (\texttt{oh},
\texttt{yd}, \texttt{eh} or \texttt{hm}) and a siman, optionally with a se'if
(\texttt{605:1}), a range of se'ifim (\texttt{658:1-2}), or a range across simanim
(\texttt{605:1-606:2}). Besides \cs{tur} and \cs{shulchanaruch} the package has
\cs{beityosef}, \cs{darkeimoshe}, \cs{bach} and \cs{aruchhashulchan}, and the commentaries
\cs{shach}, \cs{taz}, \cs{magenavraham} and \cs{mishnahberurah}, whose sub-unit is the se'if
katan. The Magen Avraham and the Mishnah Berurah cover Oraḥ Ḥayyim only, so they take just the
reference:
\begin{examples}
\row{\shulchanaruch{oh}{605:1}}
\row{\tur{yd}{114}}
\row{\shach{yd}{114:3}}
\row{\darkeimoshe{oh}{605:3}}
\row{\mishnahberurah{605:3}}
\end{examples}
\cs{declaretursefer}\verb|[options]{\command}{Hebrew}{English}| declares more such works:
\begin{Verbatim}[commandchars=+\(\)]
\declaretursefer[seifkatan]{\pitcheiteshuvah}{+texthebrew(פתחי תשובה)}{Pitchei Teshuvah}
\declaretursefer[seifkatan, section=oh]{\shaarhatziyun}
  {+texthebrew(שער הציון)}{Shaar haTziyun}
\end{Verbatim}
Its options are \texttt{seifkatan}, for a sub-unit of \texthebrew{ס״ק} rather than
\texthebrew{סע׳}; \texttt{subunit=}, for any other (the Darkei Moshe's is \texthebrew{אות});
and \texttt{section=}, for a work on one section only, whose command then takes just the
reference.

\section{Midrash and other works}\label{sec:sefer}

\cs{sefer}\verb|{key}{...}| cites any other work, by its key and the parts of the reference,
one argument each. \cs{sefer*} leaves out the work's name. The package knows these:
\begin{center}
\small\setlength{\tabcolsep}{4pt}%
\begin{tabular}{@{}lll@{}}
\toprule
Key & Arguments & Example \\
\midrule
\texttt{bereshitrabbah}, \ldots & chapter:section & \sefer{bereshitrabbah}{38:13} \\
\texttt{tanchuma} & parashah, section & \sefer{tanchuma}{noach}{1} \\
\texttt{sifra} & parashah, chapter:section & \sefer{sifra}{emor}{13:1} \\
\texttt{sifreidevarim} & piska & \sefer{sifreidevarim}{32} \\
\texttt{mekhilta} & parashah, masekhta, parashah & \sefer{mekhilta}{yitro}{BaḤodesh}{7} \\
\texttt{pirkeiderabbieliezer} & chapter & \sefer{pirkeiderabbieliezer}{35} \\
\texttt{midrashtehillim} & psalm:section & \sefer{midrashtehillim}{23:1} \\
\texttt{midrashaggadah} & book, chapter & \sefer{midrashaggadah}{gen}{18} \\
\texttt{yalkutshimoni} & book, remez & \sefer{yalkutshimoni}{gen}{110} \\
\texttt{igrotmoshe} & section, volume, siman & \sefer{igrotmoshe}{yd}{1}{62} \\
\midrule
\texttt{mishnah}, \texttt{yerushalmi} & tractate, chapter:mishnah & \sefer{mishnah}{sukkah}{4:1} \\
\texttt{rambam} & laws, chapter:law & \sefer{rambam}{teshuvah}{3:4} \\
\bottomrule
\end{tabular}
\end{center}
The ten Rabbot are \texttt{bereshitrabbah}, \texttt{shemotrabbah}, \texttt{vayikrarabbah},
\texttt{bamidbarrabbah}, \texttt{devarimrabbah}, \texttt{shirhashirimrabbah},
\texttt{ruthrabbah}, \texttt{eichahrabbah}, \texttt{koheletrabbah} and \texttt{estherrabbah}. The Sifrei is \texttt{sifreibamidbar} or
\texttt{sifreidevarim}.
A parashah is named by a key such as \texttt{noach}, \texttt{lechlecha} or
\texttt{vezothaberachah}. Hebrew has its own forms, so that
\verb|\sefer{igrotmoshe}{yd}{1}{62}| is
\begin{center}
\texthebrew{\sefer{igrotmoshe}{yd}{1}{62}}
\end{center}
The commands \cs{mishnah}, \cs{yerushalmi} and \cs{rambam} are short for
\verb|\sefer{mishnah}|, \verb|\sefer{yerushalmi}| and \verb|\sefer{rambam}|.

\subsection{Declaring a work}

\begin{Verbatim}
\declaresefer[english={template}]{keys}{Hebrew}{English}{template}
\end{Verbatim}
declares a work, anywhere before it is used. The template is what follows the work's name,
with a placeholder for each argument of \cs{sefer}, which takes them in order:
\begin{center}
\small
\begin{tabular}{@{}llll@{}}
\toprule
Placeholder & Argument & English & Hebrew \\
\midrule
\verb|#N| & a number or range (\texttt{4}, \texttt{10-12}) & 4, \tanach*{gen}{10-12} & \texthebrew{ד׳, \tanach*{gen}{10-12}} \\
\verb|#R| & a chapter and verse, as in \cs{tanach} (\texttt{38:13}) & \tanach*{gen}{38:13} & \texthebrew{\tanach*{gen}{38:13}} \\
\verb|#D| & a daf, as in \cs{bavli} (\texttt{69b}) & 69b & \texthebrew{\bavli*{69b}} \\
\verb|#T| & a section of the Tur (\texttt{oh}) & Oraḥ Ḥayyim & \texthebrew{או״ח} \\
\verb|#B| & a book of Tanakh (\texttt{gen}) & Genesis & \texthebrew{בראשית} \\
\verb|#M| & a tractate (\texttt{shabbat}) & Shabbat & \texthebrew{שבת} \\
\verb|#P| & a parashah (\texttt{noach}) & Noaḥ & \texthebrew{נח} \\
\verb|#H| & a set of laws of the Mishneh Torah (\texttt{shofar}) & \multicolumn{2}{l}{as in \cs{rambam}} \\
\verb|#S| & any text, or \verb|{Hebrew}{English}| & \multicolumn{2}{l}{as written} \\
\bottomrule
\end{tabular}
\end{center}
The template serves both languages. When English is written differently, \texttt{english=}
gives it its own template, with the same placeholders in the same order. Spaces in a template
are printed, so it usually starts with one; a space before \verb|#N|, \verb|#R| or \verb|#D|
doesn't break across lines, so a number stays with its label. The package declares Igrot Moshe
like this:
\begin{Verbatim}[commandchars=+\(\)]
\declaresefer[english={, #T #N:#N}]{igrotmoshe, igrosmoshe}
  {+texthebrew(אגרות משה)}{Igrot Moshe}{ #T +texthebrew(חלק) #N +texthebrew(סי׳) #N}
\end{Verbatim}
Text for \verb|#S| is printed as it is written, or given as \verb|{Hebrew}{English}| for each
language, as with the Mekhilta's masekhta:
\begin{Verbatim}[commandchars=+\(\)]
\sefer{mekhilta}{yitro}{{+texthebrew(בחדש)}{BaḤodesh}}{7}
\end{Verbatim}
gives
\begin{center}
\sefer{mekhilta}{yitro}{{בחדש}{BaḤodesh}}{7}\\
\texthebrew{\sefer{mekhilta}{yitro}{{בחדש}{BaḤodesh}}{7}}
\end{center}
A placeholder other than these, or an English template whose placeholders differ from the
Hebrew one's, is an error.

\section{Commentaries}\label{sec:commentaries}

\cs{tanach} and \cs{bavli} take a commentary in their optional argument. The package knows
\texttt{rashi}, \texttt{ramban}, \texttt{ibnezra}, \texttt{rashbam}, \texttt{radak},
\texttt{sforno}, \texttt{chizkuni}, \texttt{orhachaim}, \texttt{baalhaturim},
\texttt{kliyakar}, \texttt{malbim}, \texttt{onkelos} and \texttt{yonatan}, mostly on verses,
and \texttt{tosafot}, \texttt{rif}, \texttt{rabbeinuchananel}, \texttt{rashba},
\texttt{ritva}, \texttt{ran}, \texttt{meiri} and \texttt{maharsha}, mostly on a daf; any of
them works on either. \cs{declarecommentary}\verb|{keys}{Hebrew}{English}| adds more:
\begin{Verbatim}[commandchars=+\(\)]
\declarecommentary{haamekdavar}{+texthebrew(העמק דבר)}{Ha'amek Davar}
\end{Verbatim}

The words a comment opens with go in the same optional argument: \texttt{dh=} for Hebrew,
printed after \texthebrew{ד״ה}, and \texttt{sv=} for English, printed after \textit{s.v.}
Either alone serves both languages, and words with a comma go in braces:
\begin{Verbatim}[commandchars=+\(\)]
\bavli[tosafot, dh=+texthebrew(ואם תאמר), sv=ve'im tomar]{menachot}{69b}
\end{Verbatim}
gives
\begin{center}
\bavli[tosafot, sv=ve'im tomar]{menachot}{69b}\\
\texthebrew{\bavli[tosafot, dh=ואם תאמר]{menachot}{69b}}
\end{center}

\section{Ad loc}\label{sec:sham}

\cs{sham} gives \texthebrew{\sham} or ``\sham''. With a commentary, \verb|\sham[rashi]| gives
``\sham[rashi]'' or \mbox{\texthebrew{\sham[rashi]}}. With \texttt{sham=full}, it also gives the
reference it stands for, the last one made: ``Rashi ad loc (\bavli{menachot}{69b})''.

A commentary's reference stands for its verse or daf, so after
\verb|\tanach[rashi]{gen}{18:7}|, \verb|\sham[ibnezra]| is Ibn Ezra on the same verse.
\cs{bavli*} and \cs{tanach*} don't change what \cs{sham} stands for, and within the parentheses
of a full \cs{sham} the reference isn't made again.

A document that cites a source with no reference of the package's (a responsum declared
nowhere, say) and then writes \cs{sham} would make it point at the wrong source. Packages that
divide a document into sources can prevent that (section~\ref{sec:packages}); \texttt{shiur}
does.

\section{Authors}\label{sec:authors}

\begin{Verbatim}
\declareauthor[born=..., died=..., place=...]{keys}{Hebrew}{English}
\end{Verbatim}
declares an author, and \cs{attribution}\verb|{key}| gives them: their name, their dates if
known, and where they lived.
\begin{Verbatim}[commandchars=+\(\)]
\declareauthor[born=1932, place={{+texthebrew(ירושלים)}{Jerusalem}}]{asherweiss}
  {+texthebrew(הר׳ אשר וייס)}{R. Asher Weiss}
\end{Verbatim}
\declareauthor[born=1932, place={{ירושלים}{Jerusalem}}]{asherweiss}{הר׳ אשר וייס}{R. Asher Weiss}%
gives ``\attribution{asherweiss}'' and \mbox{\texthebrew{\attribution{asherweiss}}}. Each key is
optional. A year may be AM (5692) or CE (1932): one above 3760 is taken as AM. English always
gives the CE year, and Hebrew the Hebrew one, without the thousands from 5000 to 5999
(\texthebrew{תרצ״ב}) and with them otherwise (\texthebrew{ד׳ת״ת}). An AM year is taken as
AM\,$-$\,3760, which is right from January to Elul. A year that isn't a number, such as
\verb|c. 1290|, is printed as written, or given as \verb|{Hebrew}{English}|; so is a place.

A few authors of the built-in works are declared, for \cs{attribution}:
\begin{center}
\small
\begin{tabular}{@{} >{\raggedright\arraybackslash}p{0.27\linewidth}
	>{\raggedright\arraybackslash}p{0.40\linewidth}
	>{\raggedright\arraybackslash}p{0.27\linewidth} @{}}
\toprule
Keys & English & Hebrew \\
\midrule
\texttt{rashi} & \attribution{rashi} & \texthebrew{\attribution{rashi}} \\
\texttt{ramban}, \texttt{nachmanides} & \attribution{ramban} & \texthebrew{\attribution{ramban}} \\
\texttt{rambam}, \texttt{maimonides} & \attribution{rambam} & \texthebrew{\attribution{rambam}} \\
\texttt{yaakovbenasher}, \texttt{baalhaturim} & \attribution{yaakovbenasher} & \texthebrew{\attribution{yaakovbenasher}} \\
\texttt{yosefkaro}, \texttt{karo}, \texttt{mechaber} & \attribution{karo} & \texthebrew{\attribution{karo}} \\
\texttt{rema}, \texttt{isserles} & \attribution{rema} & \texthebrew{\attribution{rema}} \\
\texttt{chafetzchaim}, \texttt{yisraelmeirkagan} & \attribution{chafetzchaim} & \texthebrew{\attribution{chafetzchaim}} \\
\texttt{yechielmichelepstein}, \texttt{epstein} & \attribution{epstein} & \texthebrew{\attribution{epstein}} \\
\texttt{moshefeinstein}, \texttt{feinstein} & \attribution{feinstein} & \texthebrew{\attribution{feinstein}} \\
\bottomrule
\end{tabular}
\end{center}
The built-in works don't name them, so citing Rashi doesn't add his dates to every heading; a
document attaches one where it wants it (below).

\cs{attribution}\verb|{who}{where}{when}|, with three arguments, gives them as written,
leaving out any that is empty: \verb|\attribution{who}{}{when}| is ``who, when''.

\subsection{Works with authors}

\cs{declaresefer}, \cs{declarecommentary} and \cs{declaretursefer} take \texttt{author=}, the key
of a declared author; \cs{declaresefer} also takes \texttt{place=} and \texttt{year=}, where and
when the work was published, which replace the author's place:
\begin{Verbatim}[commandchars=+\(\)]
\declareauthor{sperling}
  {+texthebrew(הר׳ אברהם יצחק ספרלינג)}{R. Avraham Yitz\hdot ak Sperling}
\declaresefer[author=sperling, place={{+texthebrew(לבוב)}{Lvov}}, year=5651]
  {taameihaminhagim}{+texthebrew(טעמי המנהגים)}{Ta'amei haMinhagim}{ #S}
\declarecommentary[author=rashi]{rashi}{+texthebrew(רש״י)}{Rashi}
\end{Verbatim}
A reference to such a work remembers its author, for a package that typesets sources to give
under the source's heading (section~\ref{sec:packages}). In \texttt{shiur}, a source whose
heading is \verb|\sefer{taameihaminhagim}{p. 500}| is headed
\begin{center}
\textbf{Ta'amei haMinhagim p. 500}\\
\textit{R. Avraham Yitzḥak Sperling, Lvov, 1891}
\end{center}
in English, and the same in Hebrew. A subtitle the source gives itself wins, and
\texttt{attribution=off} turns this off.

\section{Dates}\label{sec:dates}

\cs{taarich}\verb|{day}{month}{year}| gives a Hebrew date in the language around:
\begin{examples}
\row{\taarich{24}{tevet}{5779}}
\row{\taarich[shabbat]{25}{kislev}{5750}}
\row{\taarich[day=tuesday]{4}{nisan}{5704}}
\row{\taarich[full]{1}{adar2}{5784}}
\row{\taarich[civil]{2019-01-01}}
\end{examples}
Months are keys like any other (\texttt{tevet}, \texttt{teves}; \texttt{cheshvan} or
\texttt{marcheshvan}, which prints as written); a leap year's Adars are \texttt{adar1} and
\texttt{adar2} (or \texttt{adarrishon}, \texttt{adarsheni}), and \texttt{adar} alone is the Adar
of an ordinary year. The options are:
\begin{description}
\item[\texttt{shabbat}, or \texttt{day=sunday} \ldots\ \texttt{friday}] the day of the week,
before the date;
\item[\texttt{full}] the Hebrew year with its thousands (\texthebrew{ה׳תשע״ט});
\item[\texttt{civil}] the date is a civil one, \texttt{2019-01-01} or \texttt{today}, given as
its Hebrew date; \cs{taarich} then takes just that one argument.
\end{description}
The date is checked: a day the month hasn't (30 Ḥeshvan in a year where Ḥeshvan has 29 days) is
an error, and a day of the week that doesn't match the date, Adar in a leap year or Adar II in
an ordinary one is a warning. The calendar follows Reingold and Dershowitz,
\textit{Calendrical Calculations}.

\cs{shanah}\verb|{5787}| gives a Hebrew year with its thousands, \texthebrew{\shanah{5787}}, and
\cs{shanah*} without them, \texthebrew{\shanah*{5787}}; for the first hundred years of a
millennium \cs{shanah*} keeps them (\texthebrew{\shanah*{5015}}), and a round thousand is
\texthebrew{\shanah{5000}}. \cs{parashah}\verb|{vayeshev}| gives ``\parashah{vayeshev}'' or
\texthebrew{\parashah{vayeshev}}, and \cs{parashah*} the name alone.

\section{Restyling and overriding}

A later declaration replaces an earlier one. So a document, or a package of your own loaded
after \texttt{mareimekomot}, can cite a built-in work in its own style, and the last
declaration wins. A declaration that names any key of an existing work or commentary replaces
it under all its keys, so redeclaring \texttt{igrotmoshe} changes \texttt{igrosmoshe} too, and
redeclaring \texttt{nachmanides} changes \verb|\tanach[ramban]|:
\begin{Verbatim}[commandchars=+\(\)]
% "Mishnah Sukkah 4:1" rather than "M. Sukkah 4:1"
\declaresefer{mishnah}{+texthebrew(משנה)}{Mishnah}{\nobreakspace#M #R}
\declarecommentary{nachmanides}{+texthebrew(רמב״ן)}{Nachmanides}
\declaretursefer{\tur}{+texthebrew(טור)}{Arba'ah Turim}
\end{Verbatim}
While the package itself loads, two of its names whose keys match alike is an error, which
keeps the built-in lists consistent; after that, a clash is a deliberate replacement.

\section{Numbers and line breaks}

Hebrew numbers take a geresh when they are one letter and gershayim before the last letter
when they are more (\texthebrew{ח׳, ל״ח, תר״ה}); 15 and 16 are \texthebrew{ט״ו} and
\texthebrew{ט״ז}. A chapter and verse is \texthebrew{ל״ז, י״ח} in Hebrew and 37:18 in
English, and a range takes an en dash.

A reference never hyphenates. It breaks only between the words of a long name, as in the
Mishneh Torah's, and at the dash of a range. Not every font has an en dash or \textit{ḥ}
(Taamey Frank CLM has no en dash, Trebuchet MS no \textit{ḥ}), so the package checks the font
in use: a range without an en dash takes a hyphen, and a missing \textit{ḥ} is made from
\textit{h} and a dot below.

\section{For package writers}\label{sec:packages}

Three commands let a package that divides a document into units (sources, sections) keep
\cs{sham} from pointing past a unit with no reference:
\begin{description}
\item[\cs{mareimekomot@startunit}] at the start of a unit;
\item[\cs{mareimekomot@endunit}] at its end: if no reference (or \cs{sham}) was made in the unit,
	\cs{sham} is cleared;
\item[\cs{mareimekomot@forget}] clears \cs{sham} at once.
\end{description}
A cleared \cs{sham} gives plain ``ad loc'', with nothing in parentheses, even with
\texttt{sham=full}.

Three more give the author of a work cited in a heading (section~\ref{sec:authors}):
\begin{description}
\item[\cs{mareimekomot@clearattribution}] before the heading;
\item[\cs{mareimekomot@ifattribution}\texttt{\{true\}\{false\}}] after it: whether a reference
	in it named an author, and \texttt{attribution} isn't \texttt{off};
\item[\cs{mareimekomot@attribution}] gives that author, in the language around.
\end{description}

\clearpage
\appendix
\section{Lists of names}

The tables below are read from \texttt{mareimekomot.sty} when this manual is compiled, so they
are always the package's own. Any key in a row names that row.

\subsection{Books of Tanakh}
\namelist{book}

\subsection{Tractates}
For \cs{bavli}, \cs{mishnah} and \cs{yerushalmi}, and \verb|#M|. The second part has the
tractates with Mishnah but no Bavli.
\namelist{tractate}

\subsection{Sections of the Tur}
For the Tur and the works arranged by it, and \verb|#T|.
\namelist{section}

\subsection{Laws of the Mishneh Torah}
For \cs{rambam} and \verb|#H|, in order.
\namelist{hilkhot}

\subsection{Parashot}
For \verb|#P|.
\namelist{parashah}

\subsection{Commentaries}
For the optional argument of \cs{tanach} and \cs{bavli}, and \cs{sham}.
\namelist{commentary}

\subsection{Months}
For \cs{taarich}. Plain \texttt{adar} is the Adar of an ordinary year.
\namelist{month}

\section{Licence}

Copyright 2023--2026 Gabriel Gendler Yom-Tov. This work may be distributed and/or modified under
the conditions of the \LaTeX{} Project Public License, either version 1.3c of this license or (at
your option) any later version. The latest version of this license is in
\url{https://www.latex-project.org/lppl.txt}.

\end{document}
